\documentclass[11pt]{article}
\usepackage[a4paper,margin=24mm]{geometry}
\usepackage{amsmath,amssymb}
\usepackage[T1]{fontenc}
\usepackage[hidelinks]{hyperref}
\newcommand{\code}[1]{\nolinkurl{#1}}
\setlength{\emergencystretch}{3em}

\title{R491: Outer Bridges III}
\author{TFPT RH programme}
\date{31 August 2026}

\begin{document}
\maketitle

\begin{center}
\fbox{\parbox{0.91\linewidth}{
\textbf{Verdict: HIDDEN GAP FOUND / DENSITY BRIDGE REDUCED.}
The requested \(8\to7\) census reduction is not attained.  This note
records the corrected target and does not make an RH claim.
}}
\end{center}

\section{Why the advertised two-piece closure is not valid}

The r489 Lean definition is
\[
\begin{gathered}
\text{\ttfamily FullWeilDyadicSampleConvergence}\\
\wedge\ \text{\ttfamily FullWeilArchContinuity}\\
\wedge\ \text{\ttfamily GridPoleHatIntegralIdentity}.
\end{gathered}
\]
Thus the completion has three components, not two.  In particular the
arch component cannot be omitted.  Before r491,
\code{fullWeilArchSide} was opaque, so no theorem could derive its
limit from the topology regardless of the quality of the grid
approximation.  R491 replaces it by the concrete r475 regularized
\(u\)-space pairing.

There were two further under-specifications.  The \(L^2\) witness had
no relation between its support diameter and
\code{supportRadius}; and uniform plus \(L^1\) convergence alone does
not control a weight asymptotic to \(1/u\) at zero.  The witness now
has a support interval of length \(R\) and a Lipschitz constant.

\section{Polar sign repair}

The r376 potential and pairing are
\[
\Pi(t)=-8(\cosh(|t|/2)-1),\qquad
P_D(g)=-\sum_k g(kD)\frac{\Delta_D^2\Pi(kD)}D .
\]
Since \(\Pi''(t)=-2\cosh(t/2)\), the outer minus gives
\[
P(g)=\int_{\mathbb R} 2\cosh(u/2)g(u)\,du ,
\]
not the negative weight written in r489.  For the one-cell test
\(D=1,x_0=1\), both the discrete pairing and the positive integral
equal \(2.04201544330209\); the negative integral has the opposite
sign.  Lean's \code{fullWeilPoleWeight} is corrected accordingly.

\section{Explicit dyadic construction}

R491 defines the represented half-open step function
\code{GridElement.toStepFun} and the actual sampled sequence
\[
 D_m=2^{-m},\quad N_m=\lfloor R2^m\rfloor,\quad
 x_i^{(m)}=h(a+iD_m).
\]
This is \code{dyadicSampleGrid h a R m}.  The values are arbitrary
reals: \code{GridElement.x : Fin steps -> Real} has no value
quantization.  Lean proves
\[
N_mD_m\le R
\]
as \code{dyadicSampleGrid_supportBound_le}.

The exact remaining convergence target is
\code{FullWeilDyadicSampleConvergence}.  It asks this explicit sequence
to satisfy the r489 fixed-support uniform-plus-\(L^1\) topology.
Its implication to the former existential density proposition is
proved.

Mathlib's \code{MemLp.exists_simpleFunc_eLpNorm_sub_lt}
approximates by general measurable simple functions.  It does not
specialize their level sets to dyadic intervals.  The explicit route
must therefore prove the sampled-step \(L^2\) estimate from the
Lipschitz/support contracts and then the autocorrelation estimate
\[
\|h_m\star\widetilde h_m-h\star\widetilde h\|_\infty
\le(\|h_m\|_2+\|h\|_2)\|h_m-h\|_2.
\]

\section{Exact pole and arch remainders}

The pole seam is narrowed to
\code{GridPoleHatIntegralIdentity}.  It says that every finite
piecewise-linear ACF equals its hat expansion and that integrating
each hat against \(-\Pi''=2\cosh(u/2)\) produces the second difference
in \code{polePairingZ}.  The implication from this pointwise identity
to the sequence-level pole dictionary is proved.

For the arch channel, a valid proof needs a common Dini/Lipschitz
majorant near \(u=0\), plus invariance of the regularized formula when
the integration endpoint is enlarged beyond the actual grid support.
Uniform plus \(L^1\) convergence by itself is insufficient for the
singular kernel.

\section{Census and remaining outer bridges}

The direct Lean census remains eight, with three sorries in
\code{ExternalBridges.lean}: the corrected completion package,
standard explicit-formula identification, and off-critical
separation.  Hence the density bridge is reduced but not closed.

After its eventual closure, bridge 2 still needs the full
Guinand--Weil arithmetic/spectral dictionary and zero-multiplicity
infrastructure absent from Mathlib.  Bridge 3 still needs the
off-critical separating smooth autocorrelation; the logical
conversion to Mathlib's \code{RiemannHypothesis} is already proved.

\end{document}
